\section{The Evil Skill Generalization Benchmark}
\label{sec:benchmark}

Evaluating self-healing and skill generalization requires testing conditions where skills \emph{will} break, and measuring whether the system can recover. Standard benchmarks (WebArena, Mind2Web) evaluate execution on canonical web pages but do not systematically stress-test the failure-and-recovery path. We propose the \emph{Evil Skill Generalization Benchmark}, a deliberately adversarial test suite designed to fill this gap.

\subsection{Design Principles}

The benchmark is built around three principles: (1) skills should break on parameter variants to test healing; (2) pages should be intentionally obfuscated to test robustness; and (3) the same page template should yield different experiences across variants to test generalization.

\subsection{Task Structure}

The benchmark comprises 60 tasks across 20 web application templates, each with 3 parameter variants. Templates span four domains: \textbf{finance} (Yahoo Finance, Google Finance, CoinMarketCap, SEC EDGAR), \textbf{e-commerce} (Amazon, Newegg, Zillow, Craigslist, eBay), \textbf{knowledge} (Wikipedia, IMDb, Goodreads, Crunchbase), and \textbf{utilities} (weather, currency, flights, packages, WHOIS, GitHub, government databases). Each template specifies expected extraction keys and a golden-path navigation sequence.

\subsection{Obfuscation Layers}

Each template page implements 7 anti-LLM obfuscation techniques:

\begin{enumerate}[leftmargin=*,itemsep=1pt]
    \item \textbf{Dynamic injection:} Navigation buttons injected by JavaScript after a deliberate delay, absent from the initial DOM.
    \item \textbf{Visibility attacks:} Interactive elements rendered with \texttt{opacity: 0.005} or \texttt{visibility: hidden}, requiring DOM-level inspection.
    \item \textbf{Multi-layer consent:} Stacked consent walls (ToS, checkbox, CAPTCHA, accordion), each layer parameterized per variant.
    \item \textbf{Ghost decoys:} 10 near-identical transparent buttons; 9 reset progress on click, only 1 advances the task.
    \item \textbf{CAPTCHA challenges:} Codes rendered on canvas elements, requiring vision or DOM-level verification.
    \item \textbf{Conditional DOM injection:} Content dynamically created after accordion expansion, not present in static HTML.
    \item \textbf{Variant-specific identifiers:} HTML element IDs incorporate variant parameters, preventing hardcoded selectors.
\end{enumerate}

\subsection{Evaluation Protocol}

Each task is evaluated in two modes: \textbf{without-skill} (baseline, the agent reasons from scratch) and \textbf{with-skill} (a skill recorded on one variant is replayed on the remaining two). Metrics include binary success/failure, partial-match scoring (exact: 1.0, partial: 0.5, normalized per-key), steps taken, wall-clock time, and token consumption. A task is considered successful if aggregate score $\geq$ 0.67.

\subsection{Implementation}

A mock web server serves the 20 obfuscated templates locally, enabling controlled, reproducible evaluation without dependency on live websites. The benchmark infrastructure includes a canonical skill recorder, a batch runner with isolated browser profiles, a results evaluation script, and a report generator. Full-scale evaluation is ongoing and will be reported in future work.
